\begin{abstract}\noindent Token-level credit assignment for language-model
reinforcement learning is usually formulated as if the policy were fully
trainable, while practical LLM-RL pipelines often rely on
parameter-efficient fine-tuning, especially LoRA. We argue that this
separation hides a structural failure mode. Under LoRA, the policy is
restricted to a low-rank neighborhood of the reference model, so the
per-token output-distribution differences used by common intrinsic credit
signals, surprisal, entropy reduction, and policy divergence, can become
degenerate after within-trajectory normalization, either approaching uniform
weights or concentrating on a small set of task-agnostic positions. We
formalize this behavior and propose measuring it directly with concentration
diagnostics such as weight Gini and effective-token ratio. We then introduce \emph{Adapter-Residual Credit
Assignment} (ARCA), a lightweight alternative that derives token salience from
the adapter's own hidden-state residual,
$\|h^{\text{adapted}}_t - h^{\text{base}}_t\|_2$. ARCA asks where the adapter
actually changes the model, rather than where the output distribution appears
uncertain or shifted, and requires no learned reward model, value head, or tree
construction. In a compact MATH/Qwen3-1.7B GRPO sweep, ARCA exhibits the
predicted non-degenerate middle-regime credit distribution under matched
rollout budgets and remains competitive with rank-matched baselines.
\end{abstract}
